\section{The Hair of the Dog}

They say that the cure for a hangover is the ``hair of the dog that bit you''. This is a homeopathic remedy in the sense that a very dilute portion of the poison that caused the illness provides the method for its cure.

Nevertheless, within the past few weeks I have heard on the Internet two serious proposals from ``defenders of Western civilization'' (albeit in quite different ways) about how to reach ``young people''. And they both proposed (independently) that what ``young people'' want consists of multimedia presentations with music, snappy graphics, flashing images, lively web sites ... Unfortunately, that is the very disease they suffer from! Perhaps there is a proper homeopathic dose to apply, but a full dose of the same poison is sure not to have the desired effect. In the 60s, Marshall McLuhan pointed out that the ``medium is the message''. Therefore, an attempt to change the content while retaining the same medium will not have the desired effect of changing the ``message''. At best, there will result isolated or shallow ``conversions'', but those loyalties can easily shift to the next big media event. And scattered counter-cultural elements will not be able to compete with the barrage of large-scale public media events.

The poison is the very dependence on sensory stimulation and its kaleidoscopic variations. The only effective antidote, then, is a shift from the stimulation of the senses to the stimulation provided by the apprehension of ideas. For this, Good and True ideas must be communicated. And their dissemination can be slow and tedious. There is no expectation that the ``masses'' will respond to such an appeal, but an elite will gradually form around them. This elite will serve as the fulcrum of the lever that will effect permanent change.

Update 24 April 2022: The website is defunct and the creator is disgraced.


\flrightit{Posted on 2010-04-24 by Cologero}

\begin{center}* * *\end{center}

\begin{footnotesize}
\begin{sffamily}
\texttt{Comprehensor on 2010-04-25 at 09:37 said:}

People young and old are surely attached to sensory phenomena.

Do you think that another part of the problem is their unwillingness to sacrifice a portion of their individuality in the pursuit of truth and unity (as opposed to individual belief and small competing factions)?

\hfill

\texttt{Cologero on 2010-04-26 at 22:04 said:}

This is not a ``put down'' of anyone attached to sensory phenomena. Merely a recognition that by the very nature of things most think only in terms of opinion. A smaller number will advance to rational thinking. But true Wisdom --- episteme --- is reserved to the few willing and able to make the necessary efforts.

Therefore, all attempts to appeal only to the sensory will be fated to partial results, which will be fickle, moving to the next ``exciting thing''. Efforts must be made to create the elite and the rest will follow.

\hfill

\texttt{Comprehensor on 2010-04-27 at 08:57 said:}

At the end of the day, tradition will only attract a select few, who more than likely will only pick and choose what they like to believe as conforms to their own individualities.

\end{sffamily}
\end{footnotesize}

